% called by main.tex
%
\chapter{Wakefields and Impedances}
\label{ch::chapter2}

\begin{center}
  \begin{minipage}{0.5\textwidth}
    \begin{small}
      In which the theoretical introduction and state of the art on electromagnetic calculation of wakefields and impedances will be detailed
    \end{small}
  \end{minipage}
  \vspace{0.5cm}
\end{center}

\section{Brief introduction to high energy particle accelerators}

\section{The concept of beam-coupling impedance}

\subsection{Resistive wall impedance}

\subsection{Direct and indirect space charge impedance}

\subsection{Impedance resonator formalism} % Maybe move to ANNEX

\subsection{Analytical impedance for simple geometries}

\section{Importance of beam-coupling impedance in particle accelerators}

\section{Impact on beam-dynamics and beam stability}
% Collective effects

\subsection{Single-bunch instabilities}

\subsection{Multi-bunch instabilities}

\section{Beam-induced heating of individual accelerator components}

\subsection{Beam power spectrum}

\subsection{Single-beam power loss}

\subsection{Two counter-rotating beams power loss}

% Currently chapter 3
%\section{How to compute beam-coupling impedance of accelerator components?}

% \subsection{Analytical calculations for simple geometries}

% \subsection{Numerical computations in frequency-domain}

% \subsection{Numerical computations in time-domain}


